
\section{Extended Related Work}
\label{sec:extended-related-work}

Here we offer an extended discussion of related work previously summarized in Section~\ref{sec:related-work}.

\subsection{Theoretical Foundations} Our work is closely related to the theoretical notion of \emph{universal induction}~\citep{solomonoff1964formal,hutter2000theory,lattimore2013no,achille2021information} -- and related resource-bounded variants~\citep{levin1973universal,nakkiran2021turing}. Previous work has developed theoretical frameworks applying these notions to density estimation~\citep{barron1985logically,barron1991minimum}, sequential decision theory~\citep{hutter2005universal}, and kernel methods~\citep{hamzi2024bridging}.  For neural networks, \citet{schmidhuber1997discovering} introduced a probabilistic search algorithm for discovering neural networks with low Kolmogorov complexity. However, none of these theoretical frameworks directly lead to practical and scalable training objectives for neural networks.

\subsection{Variational Inference}
 Variational inference was proposed as a regularizer to improve generalization~\citep{hinton1993keeping}, under the ``bits back'' argument, which can also be viewed as a form of Bayesian inference~\citep{honkela2004variational}. 
Such variational methods were popularized by the success of variational autoencoders (VAEs)~\citep{kingma2014auto}, and related advances that enabled applying such methods to neural networks using standard gradient-based optimizers~\citep{graves2011practical,blundell2015weight}. Variational methods have found success for network compression -- enabling weight pruning~\citep{louizos2017bayesian} and quantization~\citep{achterhold2018variational,ullrich2017soft} -- as well as for probing~\citep{voita2020information} and improving uncertainty modeling~\citep{sankararaman2022bayesformer,gal2016dropout}. Other work has pursued applying variational bottlenecks to network \emph{activations}~\citep{fehr2024nonparametric}, as opposed to network \emph{weights}.
However, in general, the effectiveness of variational methods for implementing MDL-inspired regularizers for improving generalization has remained elusive~\citep{blier2018description,cinquin2021pathologies}. Prior work has typically found tighter compression bounds via prequential coding methods~\citep{blier2018description,bornschein2023sequential}. Our work provides some new perspectives on the poor performance of variational approaches: objectives evaluated by prior work may have poor asymptotic bounds, and variational objectives can be difficult to optimize. 

\subsection{Non-Variational Complexity Measures}\label{sec:non-variational-related-work} As regularization is a foundational concept in machine learning, there are many related methods for neural networks~\citep{jiang2020fantastic}, and we give only a brief survey here. We focus on recently proposed compression-based methods with connections to the MDL principle, as these are most relevant to the description length objectives discussed in this paper.

Motivated by prior work on the intrinsic dimensionality of neural networks~\citep{li2018measuring}, \citet{lotfi2022pac} propose a complexity measure based on subspace sampling combined with dynamic quantization, similar to the quantization method of~\citet{han2016deep}. 
\citet{lotfi2024non} combines the method of \citet{lotfi2022pac} with LoRA~\citep{hu2022lora,dettmers2023qlora}, to further drive compression. Both \citet{lotfi2022pac} and \citet{lotfi2024non} discuss generalization bounds for neural networks related to compression, a relevant topic, but not one we study explicitly in this paper.
\citet{demoss2025complexity} studies the connection between the complexity of a neural network and ``Grokking'', i.e. the transition from memorization to generalization during training. To this end, they propose a new complexity measure and regularizer based on coarse-grained quantization and spectral entropy, which encourages low rank parameter matrices.
\citet{abudy2025minimum} similarly advocate for regularizers grounded in MDL instead of standard norm-based penalties on weights. They demonstrate how norm-based regularizers actively push RNN weights away from perfect initializations on algorithmic tasks, while their MDL-inspired objective does not. Their proposed MDL measure can be interpreted as a two-part code, using a non-differentiable prior that encodes individual weight values according to their representation as a rational number.
\citet{dwivedi2023revisiting} similarly studies MDL-inspired complexity measures in over-parameterized models where parameter count does not provide a suitable measure of complexity, and proposes a principled measure, MDL-COMP. However, their scope is limited to specific classes of linear models and kernel methods.

The variational code proposed in Section~\ref{sec:transformer-practical-code} drives compression by encouraging soft quantization of model weights around the means of the components of the GMM prior, and by reducing the effective number of weights by encouraging higher uncertainty posterior distributions close to the prior for ``unused'' capacity. The non-variational approaches discussed here similarly combine some form of quantization with some alternative means of reducing the effective number of parameters, e.g. through low-rank approximation. While prior work has not established asymptotic guarantees for these alternative encoding schemes, future work could aim to construct families of asymptotically optimal, or quasi-optimal, codes based on these alternative methods.

\subsection{Computational Universality of Transformers}
\label{sec:zmap-related-work}

There has been considerable interest in establishing the theoretical expressivity of various classes of Transformers~\citep{perez2021attention,chiang2023tighter,yun2020are,feng2023towards,merrill2024expressive,nowak2024representational,merrill2024little}. Most relevant to our constructed mapping between computable, rational-valued distributions and Transformer weights is prior work establishing various equivalences between Transformers and Turing machines. \citet{perez2021attention} established Turing completeness of an encoder-decoder Transformer in a non-probabilistic language recognition setting, with \citet{merrill2024expressive} and \citet{feng2023towards} extending this result to decoder-only variants. Most relevant to our result is \citet{nowak2024representational}, which demonstrates Turing completeness in a probabilistic setting. One difference is that our result holds for Transformer encoders (in the limit as context size and number of layers increase) compared with \citet{nowak2024representational}'s result for Transformer decoders with intermediate decoding steps. Their result also does not directly extend to prefix Turing machines, which is necessary to support our theoretical results. However, their work could be potentially helpful for future work seeking to establish the existence of families of asymptotically optimal codes for Transformer decoders with intermediate decoding steps.
Finally, \citet{schuurmans2023memory} shows that a Transformer decoder augmented with an external memory tape is computationally universal. Broadly, considering description length objectives over models with external tools could be of interest for future work. 

We use the ALTA compiler~\citep{shaw2024alta} to support our demonstration of universality. ALTA was inspired by RASP~\citep{weiss2021thinking}, an alternative language for compiling programs to Transformers, and the related Tracr~\citep{lindner2024tracr} compiler. ALTA offers better support for programs with loops, which is useful for emulating a Turing machine.

\subsection{Simplicity Bias} While our work focuses on implementing a simplicity bias \emph{explicitly} in a training objective, several studies have evaluated the \emph{implicit} simplicity bias of neural networks~\citep{zhou2023algorithms,abbe2023generalization,bhattamishra2023simplicity,tsoy2024simplicity,chen2024sudden,mingard2025deep}, or lack thereof~\cite[e.g.,][]{nikankin2025arithmetic}. Others study the simplicity bias of \emph{in-context learning} with pre-trained models~\citep{goldblum2023no,deletang2024language}, or proposed tasks to enhance this bias~\citep{grau-moya2024learning}.
